\subsection{Cubic incidence and Lorentzian realization}

Cubic completion supplies a relation between interior and boundary.
Its geometry can be developed through incidence of the six faces with
the eight vertices. Label faces by \((i,\varepsilon)\),
\(i\in\{1,2,3\}\), \(\varepsilon\in\{+1,-1\}\), and vertices by
\(\nu\in\{+1,-1\}^3\). The incidence matrix is
\[
 M_{(i,\varepsilon),\nu}=\mathbf1_{\{\nu_i=\varepsilon\}}.
\]
Let \(u=(1,\ldots,1)^{\mathsf T}\in\RR^6\), let \(R_F\) denote
face opposition, and let \(P_u=uu^{\mathsf T}/6\), \(P_-=(I-R_F)/2\).

\begin{proposition}[Four-dimensional incidence quotient]
The incidence Gram matrix satisfies
\[
 MM^{\mathsf T}=12P_u+4P_-.
\]
Its eigenvalues \(12,4,0\) have multiplicities \(1,3,2\).
Consequently,
\[
 Q_\square=\RR^6/\ker(MM^{\mathsf T})
 \simeq\RR u\oplus E_-,\qquad E_-=\operatorname{im}P_-.
\]
\end{proposition}
\begin{proof}
A face contains four vertices; two opposite faces have empty intersection,
and two distinct non-opposite faces share two vertices.
The matrix \(12P_u+4P_-\) has exactly these entries.
The uniform subspace has dimension one; differences between opposite
faces generate a three-dimensional space orthogonal to it.
Its complement has dimension two and belongs to the kernel.
\end{proof}

The decomposition \(1+3\) comes from incidence. The temporal choice
is declared through the bilinear form
\[
 B_\square([x],[y])=\langle x,(P_--P_u)y\rangle.
\]
It is well defined on the quotient because both projectors annihilate
the kernel. In the basis
\[
 t=[u]/\sqrt6,\qquad
 d_i=[e_{i,+}-e_{i,-}]/\sqrt2,
\]
its matrix is \(\operatorname{diag}(-1,1,1,1)\). The domain,
decomposition, and sign convention producing the Minkowski realization
are thus explicit. The positive Gram matrix determines the quotient;
the temporal assignment determines the indefinite form on it.

The operator
\[
 W_\square=\sqrt6P_u+\sqrt2P_-,\qquad
 W_\square e_{i,\varepsilon}=t+\varepsilon d_i
\]
sends the six face labels to null directions, since
\(B_\square(t+\varepsilon d_i,t+\varepsilon d_i)=-1+1=0\).
Direction reversal and face opposition thus acquire a concrete
geometric realization.

\subsubsection{Soldering, reversal, and subdivision}

The cellular realization uses an invertible transport \(U_m(a)\),
preserving the Lorentzian form, between the endpoint fibers of an
oriented dual edge \(a\), a face label \(\lambda_m(a)\),
and an orientation \(\sigma_m(a)\).
With \(\ell_m=\ell_0/9^m\), the soldering map is
\[
 \beta_m(a)=\ell_m\sigma_m(a)
 W_{\square,s(a)}e_{\lambda_m(a)}.
\]
Compatible reversal satisfies
\[
 \beta_m(\bar a)=-U_m(a)\beta_m(a).
\]
Identify the source fibers at two consecutive levels through the declared
refinement isometry. If the nine subsegments \(a_1,\ldots,a_9\),
at level \(m+1\), preserve label and orientation when transported
to that common fiber, then
\begin{equation}
 \beta_m(a)=\sum_{j=1}^9
 U_{m+1}(a_1\cdots a_{j-1})^{-1}\beta_{m+1}(a_j).
 \label{eq:soldadura-refinamiento}
\end{equation}
Each transported term equals \(\ell_m/9\) times the same null
direction, proving the identity. The compatibility hypothesis determines
which subdivisions realize this refinement. The equation brings together
direction, scale, and transport, rather than limiting geometric
correspondence to a dimension count.

With orientation and the Lorentzian form also fixed, Hodge duality
on \(\Lambda^2Q_\square^*\) satisfies \(\star^2=-I\).
In dimension four, the sign is \((-1)^{2(4-2)+1}=-1\).
This yields a real complex structure on two-forms and, after
complexification, the eigenspaces with eigenvalues \(+i\) and \(-i\).
The electromagnetic application of this domain belongs to Article III of
the series and uses the electronic construction of Article I together
with its stated coupling.

\subsection{Compatibility of modular and positional reductions}

The block-decimal base enters through Euclidean division preserving
residue and quotient simultaneously:
\[
 n=1000q+r,\qquad 0\le r<1000.
\]
Since \(1000\equiv1\pmod9\) and \(1000\equiv1\pmod3\),
\[
 n\equiv q+r\pmod9,\qquad n\equiv q+r\pmod3.
\]
Reduction of a block therefore requires the quotient's contribution
to recover the integer's class. The numbers \(0\) and \(1000\)
have the same remainder modulo \(1000\), but different classes
modulo \(9\) and \(3\). This is an operational reason for preserving
memory through a change of resolution.

Ternary refinement and determination of a decimal prefix can be
coordinated through a single integer remainder. If \(N/3^L\)
is a cylinder endpoint and \(D/1000^r\) a decimal endpoint, define
\[
 A=1000^rN-D3^L.
\]
On adding six ternary symbols, \(N'=729N+\nu\),
\(L'=L+6\), with \(0\le\nu<729\). If the decimal prefix remains
fixed, one obtains
\[
 A'=729A+\nu1000^r.
\]
If a decimal block \(\delta\) is also incorporated, then
\(D'=1000D+\delta\), \(r'=r+1\), and
\begin{equation}
 A'=729000A+\nu1000^{r+1}-729\delta3^L.
 \label{eq:residuo-dos-bases}
\end{equation}
Both identities follow by substitution of the definitions. Comparing
endpoints is thus reduced to integer operations retaining each
prolongation's contribution. This compatibility makes it possible
to refine a region before evaluating its real coordinate.
